%! Exponential Function Manipulation — Unit 2, Lesson 2.4 — Group Activity
\documentclass[10pt]{article}
\usepackage{precalculus-article}
\usepackage{precalculus-boxes}
\newcommand{\TallMath}[1]{$\displaystyle #1\rule[-1.4em]{0pt}{3.2em}$}

\begin{document}

\pageheader{Unit 2, Lesson 2.4}{Group Activity}
\namepartnerperiod

\vspace{0.1in}

\begin{scenariobox}[Equivalent Forms of Exponential Functions]{sky}
\small
Every exponential function can be written in many equivalent forms. Choosing the right form
makes certain features of the function immediately visible. In this activity you will practice
rewriting exponential expressions using the product property, the power property, and the
rules for negative and fractional exponents.
\end{scenariobox}

\vspace{0.1in}

\begin{tcolorbox}[colback=white, colframe=black!40,
    title=\textbf{Tier R --- Remediate},
    fonttitle=\bfseries, arc=2mm, left=3mm, right=3mm, top=2mm, bottom=2mm]

For each expression, (i) identify which property to apply and (ii) rewrite in the form $ab^x$.
Show each step.

\renewcommand{\arraystretch}{2.4}
\begin{tabularx}{\linewidth}{@{} >{\centering\arraybackslash}p{2.4cm} >{\centering\arraybackslash}p{3.2cm} X X @{}}
\textbf{Expression} & \textbf{Property} & \textbf{Step 1} & \textbf{Final: $ab^x$} \\
\midrule
$2^{x+4}$  & \blank{2.8cm} & $2^x \cdot 2^{\blank{0.7cm}}$ & \blank{3cm} \\
$3^{2x}$   & \blank{2.8cm} & $(3^{\blank{0.7cm}})^x$ & \blank{3cm} \\
$5^{-x}$   & \blank{2.8cm} & $\left(\dfrac{1}{\blank{0.8cm}}\right)^x$ & \blank{3cm} \\
$4^{x/2}$  & \blank{2.8cm} & $(4^{\blank{0.7cm}})^x$ & \blank{3cm} \\
\end{tabularx}
\end{tcolorbox}

\vspace{0.1in}

\begin{tcolorbox}[colback=white, colframe=black!40,
    title=\textbf{Tier A --- Apply},
    fonttitle=\bfseries, arc=2mm, left=3mm, right=3mm, top=2mm, bottom=2mm]

\begin{enumerate}[itemsep=12pt]

  \item Three functions are given below. Show that all three are equivalent by rewriting each
        in the form $a \cdot b^x$.

        \[
          f(x) = 4^{x+1}
          \qquad
          g(x) = 4 \cdot 4^x
          \qquad
          h(x) = 16 \cdot \left(\tfrac{1}{4}\right)^{-x+1}
        \]

        $f(x) = $ \blank{5cm} \quad $g(x) = $ \blank{4cm} (already in form) \quad $h(x) = $ \blank{5cm}

        \vspace{0.04in}
        All three simplify to: \blank{5cm}

  \item Rewrite $f(x) = 8^{x/3}$ as a single base raised to the $x$ power. Show your work.

        \vspace{0.04in}
        $8^{x/3} = (8^{\blank{0.8cm}})^x = \blank{2cm}^x = $ \blank{3cm}

  \item Rewrite $g(x) = 27 \cdot 3^{-x}$ in the form $a \cdot b^x$. Classify as growth or decay
        and state the $y$-intercept.

        \vspace{0.04in}
        $g(x) = 27 \cdot 3^{-x} = 27 \cdot \left(\blank{2cm}\right)^x$

        \vspace{0.04in}
        Growth / Decay (circle one): \blank{2cm} \quad $y$-intercept: \blank{2cm}

\end{enumerate}
\end{tcolorbox}

\vspace{0.1in}

\begin{tcolorbox}[colback=white, colframe=black!40,
    title=\textbf{Tier E --- Extend},
    fonttitle=\bfseries, arc=2mm, left=3mm, right=3mm, top=2mm, bottom=2mm]

\begin{enumerate}[itemsep=12pt]

  \item Starting with $f(x) = 6^{2x+1}$, produce \textbf{three} different equivalent forms.
        Label each form and show each step of your algebra.

        \vspace{0.06in}
        \textbf{Form 1} (vertical dilation $a \cdot b^x$): \blank{8cm}

        \vspace{0.06in}
        \textbf{Form 2} (single base to the $x$ power, $(b^c)^x$): \blank{7cm}

        \vspace{0.06in}
        \textbf{Form 3} (your choice --- be creative): \blank{8cm}

  \item For each context below, state which of your three forms from part (1) would be most
        useful, and explain why in one sentence.

        \vspace{0.04in}
        \textbf{Context A}: You want to quickly read the $y$-intercept of the function.

        \writeline

        \textbf{Context B}: You want to compare the growth rate to $g(x) = 36^x$.

        \writeline

        \textbf{Context C}: You want to identify the horizontal stretch/compression from
        $y = 6^x$.

        \writeline

\end{enumerate}
\end{tcolorbox}

\end{document}
